\begin{textsample}{NEG Dim 1 – global\_south\_2024\_05 – Score 1.00 – global\_south\_2024\_05/108076061.txt}  \label{ex:f1_neg_373}
U.S. again to veto UN General Assembly membership for Palestine The United States will again block attempts to grant Palestinians increased rights in the United Nations General Assembly.
The United States said that in spite of the expectations of a large majority favour in a vote to be taken in New York on Friday.
Advertisement Nate Evans, spokesman for the U.S.
Mission to the UN, indicated that the U.S. would again use its veto as it last did on April 18.
Palestine currently holds only"observer"status at the world body."Should the General Assembly adopt this resolution and refer the Palestinian membership application back to the Security Council, we expect a similar outcome to what occurred in April,"he said.
Advertisement A majority of the UN ' s 193 member states are expected to vote in favour of a resolution granting Palestine significantly extended rights to participate in the sessions of the UN General Assembly.
The draft resolution does not grant Palestine regular voting rights.
Advertisement Adoption of the resolution would also likely @ @ @ @ @ @ @ @ @ @ UN General Assembly had recognised Palestine as an observer state in 2012 in spite resistance from the United States.
Advertisement Palestine and the Vatican are the only two non-member states with observer status in the body.
The resolution, which was introduced by the United Arab Emirates but drafted by the Palestinians, has been the source of disagreements at UN headquarters in New York for weeks.
Advertisement The text states that the General Assembly has determined that the"State of Palestin.
Should be admitted to membership of the United Nations."It also recommends that the UN Security Council, which holds decisive power over UN membership,"reconsider the matter favourably."Advertisement Dpa has obtained the text of the draft resolution, although the provisions and language of the resolution could still change as negotiations continue.
The move by the 193-member UN General Assembly in New York, which comes against the backdrop of the ongoing war in the Gaza Strip, is also a reflection of international opinion @ @ @ @ @ @ @ @ @ @ diplomats believe that the resolution will easily achieve the necessary two-thirds majority of all votes cast in the General Assembly.
While, the influential United States, as well as China and Russia, fear a loss of control in the upgrading of regions whose statehood is disputed.

% matched lemmas: 
\end{textsample}
